% Ad Familiares 6.2 --- headnote
% ad-familiares-06-02, just before 20 April 45 BC, at Atticus's Ficulean estate

\textit{Cicero to A. Manlius Torquatus, written from Atticus's
estate at Ficulea (in the Sabine country a few miles north-east
of Rome) shortly before 20 April 45 BC (Perseus: \textit{in
Attici Ficuleano ante xii K. Mai.}). It is the second surviving
letter of the Torquatus consolation cluster (\textit{Fam.}~6.1,
6.3, 6.4, 6.2 in their probable order of composition; the
manuscript ordering is the editorial sequence used here).
Torquatus is still abroad, still a Pompeian under uncertain
clemency, and the war that has not yet ended on the African and
Spanish fronts holds his case in suspense.}

\textit{The substantive heart of the letter is the trichotomy of
\S2 --- arms prevail, peace is restored, or everything perishes
--- offered as an exhaustive cover of the possible futures, so
that Torquatus may have a steady disposition prepared for each.
The recourse, if the third branch is what comes, to the
prediction of Marcus Antonius the orator (Torquatus's
grandfather-in-law and Cicero's lost master, whom Cicero loved
to invoke as a man of foresight) is more than an ornament: it
puts the present catastrophe inside a Roman moral lineage, and
gives Torquatus a forebear from whom the consolation comes by
inheritance. The closing assertion --- that no man has a special
claim to grieve over what is happening to all, and that fortune
that is free from fault must be borne --- is the doctrine of
\textit{Fam.}~6.1.4 in epigrammatic form, the form it will
take across the rest of the cluster.}
